% part: intuitionistic-logic
% chapter: introduction
% section: axiomatic-derivations

\documentclass[../../../include/open-logic-section]{subfiles}

\begin{document}

\olfileid{int}{int}{axd}

\olsection{بديهي \usetoken{P}{derivation}}

د شهودي قضييز منطق بديهي !!{derivation}s د تصور له پلوه تر ټولو ساده
او په تاريخي لحاظ لومړني !!{derivation} نظامونه دي۔ د کلاسيک
قضييز منطق په څېر کار کوي۔

\begin{defn}[!!^{derivability}]
که $\Gamma$ د $\Lang L$ د !!{formula}s يو سټ وي، نو له
$\Gamma$ څخه \emph{!!{derivation}} د !!{formula}s يوه متناهي
لړۍ $!A_1$، \dots،~$!A_n$ ده چې د هر $i \le n$ لپاره لاندې
يو شرط پوره کېږي:
\begin{enumerate}
\item $!A_i \in \Gamma$ وي؛ يا
\item $!A_i$ يو بديهي اصل وي؛ يا
\item $!A_i$ د ځينو $!A_j$ او $!A_k$ څخه، چې $j < i$ او $k <
  i$ وي، د مودوس پوننس له مخې راشي؛ يعنې $!A_k \ident !A_j \lif !A_i$۔
\end{enumerate}
\end{defn}

\begin{defn}[بديهي اصول]
د شهودي قضييز منطق د \emph{بديهي اصولو} سټ~$\PAx$ د لاندې
بڼو ټول !!{formula}s دي:
\begin{align}
  & (!A \land !B) \lif !A \ollabel{ax:land1}\\
  & (!A \land !B) \lif !B \ollabel{ax:land2}\\
  & !A \lif (!B \lif (!A \land !B)) \ollabel{ax:land3}\\
  & !A \lif (!A \lor !B) \ollabel{ax:lor1}\\
  & !A \lif (!B \lor !A) \ollabel{ax:lor2}\\
  & (!A \lif !C) \lif ((!B \lif !C) \lif ((!A \lor !B) \lif !C)) \ollabel{ax:lor3}\\
  & !A \lif (!B \lif !A) \ollabel{ax:lif1}\\
  & (!A \lif (!B \lif !C)) \lif ((!A \lif !B) \lif (!A \lif !C)) \ollabel{ax:lif2}\\
  & \lfalse \lif !A \ollabel{ax:lfalse1}
\end{align}
\end{defn}

\begin{defn}[!!^{derivability}]
!!^a{formula}~$!A$ له $\Gamma$ څخه هله \emph{!!{derivable}}
ګڼل کېږي، او دا په $\Gamma \Proves !A$ ليکو، چې له
$\Gamma$ څخه داسې !!a{derivation} وي چې په $!A$ پاے ته رسېږي۔
\end{defn}

\begin{defn}[قضيې]
!!^a{formula}~$!A$ هله \emph{قضيه} ده چې له تش سټ څخه د~$!A$
!!a{derivation} وي۔ که $!A$ قضيه وي، $\Proves !A$ ليکو، او که
نه وي، $\Proves/ !A$ ليکو۔
\end{defn}

\begin{prop}
  که د شهودي منطق په بديهي !!{derivation} نظام کښې
  $\Gamma \Proves !A$ وي، نو په کلاسيک منطق کښې هم
  $\Gamma \Proves !A$ وي۔ په تېره بيا، که $!A$ شهودي قضيه وي،
  نو کلاسيکه قضيه هم ده۔
\end{prop}

\begin{proof}
  هر شهودي بديهي اصل هم کلاسيک بديهي اصل دے؛ نو په شهودي منطق کښې
  هر !!{derivation} په کلاسيک منطق کښې هم !!a{derivation} دے۔
\end{proof}

\end{document}
